
\uline{le sous-schéma fermé de} \(X\) \uline{défini par} \(\underline{J}\), \(X'\) \uline{le schéma obtenu en faisant éclater} \(Y\), \(P\) \uline{le fibré projectif associé à} \(\underline{E}\), \(p\) \uline{le morphisme structural} \(P\to X\). \uline{Sur} \(P\), \uline{on considère la suite exacte canonique} :
\[
0\to \underline{F}\to p^*(\underline{E})\to\underline{O}_P(1)\to0,
\]
\uline{et le morphisme composé} :
\[
 v:\underline{F}\longrightarrow p^*(\underline{E})\xrightarrow{p^*(u)}\underline{O}_P .
\]

\uline{Alors} :
\begin{enumerate}[label=\roman*)]
\item \uline{La formation de} \(X'\) \uline{commute à tout changement de base} \(f:T\to X\) \uline{tel que} \(f^*(u)\) \uline{soit régulier}.
\item \uline{L'immersion fermée} \(j:X'\to P\) \uline{définie par l'homomorphisme surjectif} \(\mathbf{S}(\underline{E})\longrightarrow\bigoplus_{n\ge0}\underline{J}^n\) \uline{défini par} \(u\) \uline{en degré} \(1\) \uline{est une immersion régulière}, \uline{son idéal étant l'image du morphisme régulier} \(v\) ; \uline{en particulier si} \(\operatorname{codim}(i)\) \uline{est constante}, \(\operatorname{codim}(j)=\operatorname{codim}(i)-1\).
\item \uline{Le} \(\mathbf{S}(\underline{E})\)-\uline{Module} \(\bigoplus_{n\ge0}\underline{J}^n\) \uline{est parfait en tant que} \(\mathbf{S}(\underline{E})\)-\uline{Module gradué}, \uline{i}.\uline{e}. \uline{il existe localement une résolution finie}, \uline{graduée}, \uline{de} \(\bigoplus_{n\ge0}\underline{J}^n\) \uline{par des} \(\mathbf{S}(\underline{E})\)-\uline{Modules gradués localement libres de type fini} (VI 1.2).
\end{enumerate}

i) Comme \(X'=\Proj(\bigoplus_{n\ge0}\underline{J}^n)\), le résultat est une conséquence immédiate des relations (1.2.3) et de EGA II.

ii) Les propriétés énoncées étant locales sur \(X\), on peut supposer \(X\) affine, et de plus \(\underline{E}\) libre sur \(X\). On pose \(X=\Spec(A)\).

On fait d'abord la démonstration dans le cas standard \(A=\mathbb Z[T_1,\ldots,T_d]\), \(E\) étant muni d'une base \(e_1,\ldots,e_d\), avec \(u(e_i)=T_i\). Alors \(P=\Proj\bigl(\mathbb Z[T_1,\ldots,T_d,X_1,\ldots,X_d]\bigr)\), la graduation étant prise par rapport aux \(X_i\). Pour montrer que \(v\) est un morphisme régulier dont l'image est l'Idéal de \(j\), il suffit de le faire localement sur \(P\), donc sur les
